% !TeX root = ../navier-stokes-manuscript.tex
% ============================================================
% 2.1 Vortex Stretching
% ============================================================

Axial stretching lengthens one narrow material vortex-tube segment; incompressibility contracts its cross-section, draws its rotating fluid inward, and makes the interior spin faster. This is the positive stretching contribution, while viscosity remains part of the signed evolution that decides whether the rotation strengthens overall. The velocity field carrying the segment is governed by the momentum equation
\[
  \underbrace{\partial_t u}_{\text{local acceleration}}
  +
  \underbrace{(u\cdot\nabla)u}_{\text{nonlinear transport}}
  =
  \underbrace{-\nabla p}_{\text{pressure}}
  +
  \underbrace{\nu\Delta u}_{\text{viscous diffusion}}.
\]
Taking its curl turns nonlinear transport into vorticity advection and stretching, giving the material rotational balance used below. The local picture is axisymmetric about the segment's axis, with vorticity and the stretching direction aligned; within that license, the geometry can be calculated without confusing positive stretching with net growth.

\subsection{Vortex Stretching}\label{sub:ThePerfectlyStretchingVortex}

Let \(e\) be the segment's axis and \(L(t)\) its length. If the strain is locally uniform across this narrow segment and extends it along \(e\), then
\[
  S:=\frac12\bigl(\nabla u+(\nabla u)^{\mathsf T}\bigr),
  \qquad
  Se=\sigma(t)e,
  \qquad
  \sigma(t)>0,
  \qquad
  \frac{\dot L(t)}{L(t)}
  =
  e\cdot Se
  =
  \sigma(t).
\]
The segment length grows at the positive fractional rate \(\sigma(t)\). Its material volume \(\mathcal V\) cannot change, so its effective transverse area is \(A(t):=\mathcal V/L(t)\), and its effective radius is \(r_{\mathrm{eff}}(t):=\sqrt{A(t)/\pi}\). Incompressibility then gives
\[
  \nabla\cdot u=0,
  \qquad
  \frac{d}{dt}(A L)=0,
  \qquad
  \frac{\dot A}{A}=-\sigma(t),
  \qquad
  \frac{\dot r_{\mathrm{eff}}}{r_{\mathrm{eff}}}
  =
  -\frac{\sigma(t)}{2}.
\]
As the segment lengthens, its area contracts at rate \(\sigma(t)\) and its effective radius contracts at half that rate. With \(\omega:=\nabla\times u\), the curl of the momentum equation gives
\[
  \frac{D\omega}{Dt}
  =
  S\omega+\nu\Delta\omega,
  \qquad
  \frac{D}{Dt}
  :=
  \partial_t+u\cdot\nabla,
\]
and the local alignment \(\omega=|\omega|e\) turns the stretching part of this material evolution into
\[
  S\omega
  =
  \sigma(t)\omega,
  \qquad
  \frac12\frac{D|\omega|^2}{Dt}
  =
  \sigma(t)|\omega|^2
  +
  \nu\,\omega\cdot\Delta\omega.
\]
The first term is positive and spins up the narrowing segment, but the viscous term has its own sign, so the complete right-hand side decides whether \(|\omega|^2\) grows. This local concentration is compatible with finite kinetic energy. Indeed, the identity proved in \Cref{prop:ClassicalKineticEnergyIdentity}, together with the antisymmetric part of \(\nabla u\), gives
\[
  \norm{u(t)}_2
  \leq
  \norm{u_0}_2
  <\infty,
  \qquad
  \left|\frac12\bigl(\nabla u-(\nabla u)^{\mathsf T}\bigr)\right|_{\mathrm F}
  =
  \frac{|\omega|}{\sqrt2}
  \leq
  |\nabla u|_{\mathrm F}.
\]
Thus finite global energy does not prevent the rotational part of the gradient from concentrating inside the contracting segment. That coexistence creates the need to compare the full velocity field with a solution at a finer scale.

\divider{\NavierStokes{} Scaling}

Fix \(t_0<T_*\) and write \(\ell:=r_{\mathrm{eff}}(t_0)\). For \(\lambda>1\), define the exact \NavierStokes{} rescaling
\[
  u_\lambda(x,t)
  :=
  \lambda u(\lambda x,\lambda^2t),
  \qquad
  p_\lambda(x,t)
  :=
  \lambda^2p(\lambda x,\lambda^2t).
\]
At time \(\lambda^{-2}t_0\), the corresponding segment in \(u_\lambda\) has effective radius \(\lambda^{-1}\ell\). The field \(u_\lambda\) is another solution at another scale, not a later material state of the opening segment. Its momentum terms transform as
\[
\begin{aligned}
  \partial_tu_\lambda(x,t)
  &=
  \lambda^3(\partial_tu)(\lambda x,\lambda^2t),
  \\
  (u_\lambda\cdot\nabla)u_\lambda(x,t)
  &=
  \lambda^3\bigl((u\cdot\nabla)u\bigr)(\lambda x,\lambda^2t),
  \\
  \nabla p_\lambda(x,t)
  &=
  \lambda^3(\nabla p)(\lambda x,\lambda^2t),
  \\
  \nu\Delta u_\lambda(x,t)
  &=
  \lambda^3\nu(\Delta u)(\lambda x,\lambda^2t).
\end{aligned}
\]
Every contribution acquires the same cubic factor. At finer width, viscosity gains no relative advantage against nonlinear transport, while pressure and local acceleration remain in the same balance.

\divider{Critical Height}

The full equation retains its form, but homogeneous Sobolev norms change according to their order. Let \(\Lambda:=(-\Delta)^{1/2}\). Since
\[
  \widehat{u_\lambda}(\xi,t)
  =
  \lambda^{-2}\widehat u(\lambda^{-1}\xi,\lambda^2t),
\]
a change of variables yields
\[
  \norm{u_\lambda(t)}_{\dot H^s}^2
  =
  \lambda^{2s-1}
  \norm{u(\lambda^2t)}_{\dot H^s}^2.
\]
The rescaled \(\dot H^s\) norm loses scale weight for \(s<\tfrac12\), is neutral at \(s=\tfrac12\), and gains scale weight for \(s>\tfrac12\). Define the half-derivative height by
\[
  H(t)
  :=
  \frac12\norm{u(t)}_{\dot H^{1/2}}^2
  =
  \frac12\norm{\Lambda^{1/2}u(t)}_2^2.
\]
For the rescaled solution, let \(H_\lambda(t):=\tfrac12\norm{u_\lambda(t)}_{\dot H^{1/2}}^2\). The kinetic-energy norm, critical height, and first-derivative norm satisfy
\[
  \norm{u_\lambda(t)}_2^2
  =
  \lambda^{-1}\norm{u(\lambda^2t)}_2^2,
  \qquad
  H_\lambda(t)=H(\lambda^2t),
  \qquad
  \norm{\nabla u_\lambda(t)}_2^2
  =
  \lambda\norm{\nabla u(\lambda^2t)}_2^2.
\]
The finer-scale solution has less kinetic energy and a larger first-derivative norm, while its half-derivative height is scale neutral. Neutrality fixes no sign for \(H'(t)\); that sign comes from the evolution.

\divider{Nonlinear Production versus Viscous Dissipation}

At critical height, nonlinear transport contributes the signed production
\[
  \mathcal P_{1/2}[u](t)
  :=
  -\operatorname{Re}
  \pair
    {\Lambda^{1/2}u(t)}
    {\Lambda^{1/2}\bigl((u(t)\cdot\nabla)u(t)\bigr)}_{L^2}.
\]
Because \(\nabla\cdot u=0\), pressure makes no contribution to the pairing, while viscosity dissipates \(\nu\norm{\Lambda^{3/2}u}_2^2\). Consequently,
\[
  H'(t)
  +
  \nu\norm{\Lambda^{3/2}u(t)}_2^2
  =
  \mathcal P_{1/2}[u](t).
\]
The critical height rises exactly when signed nonlinear production exceeds positive viscous dissipation. Under rescaling, the two quantities become
\[
  \mathcal P_{1/2}[u_\lambda](t)
  =
  \lambda^2\mathcal P_{1/2}[u](\lambda^2t),
  \qquad
  \nu\norm{\Lambda^{3/2}u_\lambda(t)}_2^2
  =
  \lambda^2\nu\norm{\Lambda^{3/2}u(\lambda^2t)}_2^2.
\]
Both gain the same quadratic factor, so scale gives neither production nor dissipation an advantage. The rescaling contracts time by \(\lambda^{-2}\), but it does not determine the duration of a nonlinear transition in one velocity history.

\divider{The Heat Clock}

The linear heat equation \(v_t=\nu\Delta v\) supplies a viscous comparison. In Fourier variables,
\[
  \widehat v(\xi,t+\delta t)
  =
  e^{-\nu|\xi|^2\delta t}\widehat v(\xi,t).
\]
A scale-localized velocity structure of width \(r\) occupies modes with \(|\xi|\simeq r^{-1}\). Those modes change by order one when
\[
  \nu|\xi|^2\delta t\simeq1,
\]
so their linear viscous comparison time is
\[
  \tau_\nu(r)
  :=
  \frac{r^2}{\nu}.
\]
Halving the localization width quarters this comparison time. More generally,
\[
  \tau_\nu(\lambda^{-1}r)
  =
  \lambda^{-2}\tau_\nu(r).
\]
This matches the time scaling of the full \NavierStokes{} transformation, but it is not an asserted lifetime or transition duration for a nonlinear structure.

\divider{The Zeno Cascade}

Consider dyadic localization widths and their heat comparison times. For
\[
  r_n:=2^{-n}r_0,
  \qquad
  \tau_n:=\frac{r_n^2}{\nu},
\]
the successive comparison times satisfy
\[
  \tau_n
  =
  \frac{r_0^2}{\nu}4^{-n},
  \qquad
  \sum_{n=0}^{\infty}\tau_n
  =
  \frac{4r_0^2}{3\nu}
  <
  \infty.
\]
Their finite sum gives a possible schedule, not a realized cascade. Realization would require one \NavierStokes{} velocity field to exhibit scale-localized structures at widths
\[
  r_0>r_1>r_2>\cdots\longrightarrow0
\]
at sampled times
\[
  t_0<t_1<t_2<\cdots\uparrow T_*<\infty,
  \qquad
  t_{n+1}-t_n\asymp\tau_n,
\]
with comparison constants independent of \(n\). Under that condition, the remaining time satisfies
\[
  T_*-t_n
  \asymp
  \sum_{k=n}^{\infty}\tau_k
  =
  \frac{4r_n^2}{3\nu}.
\]
At width \(r_n\), both the next sampled gap and the total time left are comparable to \(r_n^2/\nu\). These intervals collapse as \(t\uparrow T_*\), but the conditional schedule still does not control the first or successive Eulerian time derivatives of that same velocity field at one fixed spatial point.
